\documentclass[11pt]{article}

% Load the geometry package to set margins
\usepackage[lmargin=2cm,rmargin=2cm,tmargin=1.8cm,bmargin=1.8cm]{geometry}

% increase nesting depth

\usepackage{enumitem}
\setlistdepth{9}
%
\renewlist{itemize}{itemize}{9}
\setlist[itemize,1]{label=\textbullet}
\setlist[itemize,2]{label=--}
\setlist[itemize,3]{label=*}
\setlist[itemize,4]{label=•}
\setlist[itemize,5]{label=–}
\setlist[itemize,6]{label=>}
\setlist[itemize,7]{label=»}
\setlist[itemize,8]{label=›}
\setlist[itemize,9]{label=·}
%
\renewlist{enumerate}{enumerate}{9}
\setlist[enumerate,1]{label=\arabic*.,ref=\arabic*}
\setlist[enumerate,2]{label=\alph*.),ref=\theenumi\alph*}
\setlist[enumerate,3]{label=\roman*.),ref=\theenumii\roman*}
\setlist[enumerate,4]{label=\Alph*.),ref=\theenumiii\Alph*}
\setlist[enumerate,5]{label=\Roman*.),ref=\theenumiv\Roman*}
\setlist[enumerate,6]{label=\arabic*),ref=\theenumv\arabic*}
\setlist[enumerate,7]{label=\alph*),ref=\theenumvi\alph*}
\setlist[enumerate,8]{label=\roman*),ref=\theenumvii\roman*}
\setlist[enumerate,9]{label=\Alph*),ref=\theenumviii\Alph*}


% Load CM Bright for math
\usepackage{amsmath}  % Standard math package
\usepackage{amssymb}  % Additional math symbols
\usepackage{cmbright} % Sans-serif math font that complements Fira Sans

\usepackage{fourier}

% Font configuration
% \usepackage{bera}
% or
% Load Fira Sans for text
\usepackage{fontspec}
\setmainfont{Fira Sans}  % System-installed Fira Sans
\renewcommand{\familydefault}{\sfdefault}  % Set sans-serif as default

% pseudo-code with math
\usepackage{listings}
\usepackage{float}
\usepackage{xcolor}
\usepackage{colortbl}
\usepackage{booktabs}
% Set TT font
% \usepackage{inconsolata}
% or
\setmonofont{IBMPlexMono-Light}
% Define custom settings for listings
\lstset{
  language=Python,
  basicstyle=\ttfamily\small,        % Monospaced font
  commentstyle=\itshape\color{gray}, % Italic and gray for comments
  keywordstyle=\color{blue},         % Keywords in blue
  stringstyle=\color{red},           % Strings in red
  mathescape=true,                   % Enable math in comments
  breaklines=true,                   % Break long lines
  numbers=left,                      % Add line numbers
  numberstyle=\tiny\color{gray},     % Style for line numbers
  frame=single,                      % Add a frame around the code
  escapeinside={(*@}{@*)},           % Allow LaTeX commands (e.g. hyperlinks)
}

\usepackage{newfloat}  % Allows creating custom float types

% Define 'listing' as a floating environment
\DeclareFloatingEnvironment[
fileext=lol,
listname=List of Listings,
name=Listing
]{listing}

% To prevent floats from moving past a section boundary but still allow some floating:
\usepackage{placeins}
% used with \FloatBarrier

\usepackage[utf8]{inputenc}
\usepackage[T1]{fontenc}
\usepackage{graphicx}
\usepackage{longtable}
\usepackage{wrapfig}
\usepackage{rotating}
\usepackage[normalem]{ulem}
\usepackage{amsmath}
\usepackage{amssymb}
\usepackage{capt-of}
\usepackage{hyperref}
\usepackage{algorithm}
\usepackage{algpseudocode}

% Title, Author, and Date (or Report Number)
\title{MRVA for CodeQL:\\Worker Pool Architecture}
\author{Michael Hohn}
\date{Technical Report 20260311}

\hypersetup{
  pdfauthor={Michael Hohn},
  pdftitle={MRVA for CodeQL: Worker Pool Architecture},
  pdfkeywords={CodeQL, MRVA, Argo Workflows, Worker Pool, Postgres MVCC},
  pdfsubject={},
  pdfcreator={Claude Code},
  pdflang={English}}

% Tiny clickable circle linking to Go source code.
% Usage: \gomark{pkg/pool/pool.go}{120}
% Opens vscode at that file and line.
\newcommand{\reporoot}{/Users/mhohn5/c/example/psa-mrva/mrvacommander}
% \newcommand{\gomark}[2]{\href{vscode://file\reporoot/#1:#2}{\textcolor{blue!60}{$\circ$}}}
\newcommand{\gomark}[2]{\href{vscode://file\reporoot/#1:#2}{\textcolor{gray!90}{$\blacklozenge$}}}

\begin{document}

\maketitle
\tableofcontents

\section{Worker Pool Motivation}

The current Argo Workflows architecture creates a new Kubernetes pod $\omega_i$ for
each job $J_i$. This incurs container startup overhead per job: image pull, process
initialization, CodeQL CLI setup, and credential exchange. For large job batches, this
overhead dominates execution time.

The worker pool $P_\omega$ addresses this by maintaining a set of long-lived worker
pods that re-use already-running containers. Workers request new jobs upon completing
their current work instead of exiting. Key design goals:

\begin{itemize}
\item \textbf{Container re-use}: Workers persist across jobs, eliminating per-job
  startup cost.
\item \textbf{No leader election}: PostgreSQL MVCC replaces leader election for
  coordination. Workers claim jobs directly via atomic SQL operations.
\item \textbf{Adaptive scaling}: Fast ramp-up when work arrives, slow decay to zero
  workers when idle.
\item \textbf{Minimal scheduler change}: The server-side scheduler main loop remains
  essentially unchanged; the pool is an added component.
\item \textbf{All state in PostgreSQL}: Worker registration, job claiming, heartbeats,
  and scaling decisions use the database as the single source of truth.
\end{itemize}

\subsection{Architectural Change Summary}

The worker pool $P_\omega$ replaces the per-job Argo workflow submission with a
pool of persistent workers:

\begin{center}
\begin{tabular}{p{4cm}|p{5cm}|p{5cm}}
\toprule
\textbf{Aspect} & \textbf{Current (Per-Job Pods)} & \textbf{New (Worker Pool)} \\
\midrule
Worker Lifecycle & New pod per job $J_i$ & Long-lived pod, multiple jobs \\
Startup Cost & Per job ($\sim$30--60s) & Once per worker \\
Job Assignment & Scheduler $\to$ Argo submit & Worker claims from Postgres \\
Coordination & Leader-elected scheduler & Postgres MVCC \\
Scaling & Fixed $W$ max concurrent & Adaptive: $W_{min}$ to $W_{max}$ \\
Idle Behavior & Pod exits & Worker polls, exits after $T_{idle}$ \\
\midrule
Flow Notation & $\Sigma_L \xrightarrow{\text{submit}} \Omega \to \omega_i$ &
  $\omega_k \xleftarrow{\text{claim}} D$ \\
\bottomrule
\end{tabular}
\end{center}

\subsection{Component Responsibilities}

The pool $P_\omega$ touches three codebases:

\begin{center}
\begin{tabular}{p{4cm}p{4cm}p{6cm}}
\toprule
\textbf{Component} & \textbf{Repository} & \textbf{Role in Pool} \\
\midrule
Shared library & \texttt{mrvacommander/} & Pool interface types, Postgres queries for
  job claiming, worker registration, heartbeat, scaling metrics \\
Agent & \texttt{agent-commander-argo/} & Worker loop: claim $\to$ execute $\to$
  complete $\to$ claim next \\
Server & \texttt{server-commander/} & Pool manager: monitor load, scale
  workers up/down via Argo API \\
\bottomrule
\end{tabular}
\end{center}

\section{Symbols and Notation}
\label{sec:symbols-notation}

We extend the notation from the Argo Scheduler specification with pool-specific
symbols.

\begin{center}
  \begin{longtable}{p{4cm}p{2.5cm}p{7.5cm}}
    \toprule
    \textbf{Concept} & \textbf{Symbol} & \textbf{Description} \\
    \midrule
    \endfirsthead

    \toprule
    \textbf{Concept} & \textbf{Symbol} & \textbf{Description} \\
    \midrule
    \endhead

    \midrule
    \endfoot

    \bottomrule
    \endlastfoot

    \multicolumn{3}{l}{\textit{Inherited Symbols (unchanged)}} \\
    Client & $C$ & Source of query submission \\
    Server & $S$ & Central API coordinator \\
    Database & $D$ & PostgreSQL database storing all state \\
    Query Suite & $\mathcal{Q}$ & Collection of queries submitted by client \\
    Repository List & $\mathcal{R}$ & Collection of repositories \\
    Job Specification & $J_i$ & Job identifier: $(s, o, r)$ \\
    Job Status & $\sigma_i$ & Pending, Queued, InProgress, Succeeded, Failed \\
    User & $u \in U$ & Individual user identified by JWT \\
    Priority Function & $\pi(u)$ & User priority for fair scheduling \\
    Max Workers & $W$ & Maximum concurrent jobs (inherited, now $= W_{max}$) \\
    \midrule

    \multicolumn{3}{l}{\textit{Worker Pool}} \\
    Worker Pool & $P_\omega$ & Set of active worker pods \\
    Worker & $\omega_k$ & $k$-th worker pod in the pool \\
    Worker Status & $\varsigma_k$ & Idle, Busy, Draining, Exited \\
    Worker Count & $|P_\omega|$ & Number of active workers \\
    \midrule

    \multicolumn{3}{l}{\textit{Scaling Parameters}} \\
    Min Workers & $W_{min}$ & Minimum pool size (default: 0) \\
    Max Workers & $W_{max}$ & Maximum pool size (default: 10) \\
    Active Workers & $W_{active}$ & Currently running workers: $|P_\omega|$ \\
    Idle Workers & $W_{idle}$ & Workers not executing a job \\
    Ramp Factor & $\lambda$ & Scale-up multiplier (default: 2) \\
    Decay Rate & $\gamma$ & Workers removed per decay interval (default: 1) \\
    \midrule

    \multicolumn{3}{l}{\textit{Timing}} \\
    Idle Timeout & $T_{idle}$ & Time before idle worker exits (default: 300s) \\
    Heartbeat Interval & $T_{hb}$ & Worker heartbeat period (default: 10s) \\
    Heartbeat Expiry & $T_{hb}^{exp}$ & Dead worker threshold (default: 30s) \\
    Scale Check Interval & $T_{scale}$ & Pool size evaluation period (default: 15s) \\
    Decay Interval & $T_{decay}$ & Minimum time between scale-down steps (default: 60s) \\
    Claim Poll Interval & $T_{claim}$ & Worker job-claim polling period (default: 2s) \\
    \midrule

    \multicolumn{3}{l}{\textit{Fairness (embedded in SQL)}} \\
    Running Count & $n_u^r$ & Running jobs for user $u$ \\
    Pending Count & $n_u^p$ & Pending/queued jobs for user $u$ \\
    Recent Runtime & $\tau_u$ & Recent runtime for user $u$ \\
    Runtime Window & $T_w$ & Time window for $\tau_u$ (default: 3600s) \\

  \end{longtable}
\end{center}


\section{Database Schema Extensions}
\label{sec:db-schema}

\subsection{Extended \texttt{job\_info} Table}

The \texttt{job\_info} table gains a \texttt{worker\_id} column and a new
\texttt{Queued} status:

\begin{center}
\begin{longtable}{p{4.5cm}p{2.5cm}p{7cm}}
\toprule
\textbf{Column} & \textbf{Type} & \textbf{Description} (Symbol) \\
\midrule
\endfirsthead

\toprule
\textbf{Column} & \textbf{Type} & \textbf{Description} (Symbol) \\
\midrule
\endhead

\midrule
\endfoot

\bottomrule
\endlastfoot

\multicolumn{3}{l}{\textit{Existing columns (unchanged)}} \\
\texttt{session\_id} & INTEGER & Session identifier \\
\texttt{owner} & TEXT & Repository owner \\
\texttt{repo} & TEXT & Repository name \\
\texttt{payload} & JSONB & Job metadata \\
\texttt{user\_id} & TEXT & User from JWT ($u_i$) \\
\texttt{started\_at} & TIMESTAMP & Start time ($t_i^s$) \\
\texttt{completed\_at} & TIMESTAMP & Completion time ($t_i^c$) \\
\texttt{argo\_workflow\_id} & TEXT & Argo metadata ($\omega_i^{id}$) \\
\hline
\multicolumn{3}{l}{\textit{New column}} \\
\texttt{worker\_id} & TEXT & Claiming worker ($\omega_k$), NULL if unclaimed \\
\end{longtable}
\end{center}

Index addition:
\begin{itemize}
\item \texttt{idx\_job\_info\_status\_created} on \texttt{(status, created\_at)} where
  \texttt{status IN ('Pending', 'Queued')}
\end{itemize}

\subsection{New \texttt{worker\_pool} Table}

The \texttt{worker\_pool} table tracks active workers:

\begin{center}
\begin{tabular}{p{4.5cm}p{2.5cm}p{7cm}}
\toprule
\textbf{Column} & \textbf{Type} & \textbf{Description} \\
\midrule
\texttt{worker\_id} & TEXT & Worker identifier (PRIMARY KEY) \\
\texttt{status} & TEXT & Idle, Busy, Draining, Exited ($\varsigma_k$) \\
\texttt{registered\_at} & TIMESTAMP & When worker joined pool \\
\texttt{heartbeat\_at} & TIMESTAMP & Last heartbeat \\
\texttt{current\_job} & TEXT & Current job key, NULL if idle \\
\texttt{jobs\_completed} & INTEGER & Total jobs completed by this worker \\
\texttt{argo\_pod\_name} & TEXT & Kubernetes pod name \\
\bottomrule
\end{tabular}
\end{center}


\section{Job Lifecycle with Worker Pool}

Every job follows an extended state transition:

\begin{equation}
\sigma_i: \text{Pending}
  \xrightarrow{\text{scheduler}}  \text{Queued}
  \xrightarrow{\omega_k\text{ claims}} \text{InProgress}
  \to \{\text{Succeeded}, \text{Failed}\}
\end{equation}

\begin{itemize}
\item \textbf{Pending}: Server has stored the job. The scheduler has not yet
  evaluated it.
\item \textbf{Queued}: The scheduler's fair-scheduling pass has selected this job as
  eligible for execution. Workers may claim it.
\item \textbf{InProgress}: A worker $\omega_k$ has atomically claimed the job via
  Postgres \texttt{FOR UPDATE SKIP LOCKED}.
\item \textbf{Succeeded / Failed}: Terminal states, unchanged from current
  architecture.
\end{itemize}

The \texttt{Queued} status decouples fair scheduling from job execution: the scheduler
controls \emph{which} jobs become available; workers control \emph{when} they execute.


\section{Fair Scheduling via Postgres MVCC}
\label{sec:fair-scheduling}

Fair scheduling is preserved without leader election by splitting it into two
operations, both using Postgres as the coordination mechanism.

\subsection{Step 1: Server Queues Jobs (Fair Selection)}

The server (or any server replica) periodically runs the fair-scheduling algorithm and
transitions selected jobs from \texttt{Pending} to \texttt{Queued}. To prevent
concurrent server replicas from both computing available capacity from the same
snapshot and over-allocating, the procedure acquires a Postgres advisory lock.

\textbf{Why advisory lock, not leader election}: A Postgres advisory lock
(\texttt{pg\_try\_advisory\_lock}) is a single atomic call with no infrastructure
overhead -- no heartbeat table, no expiration tracking, no instance registration. It
is released automatically when the transaction (or session) ends, including on
crashes. This is qualitatively different from the leader election protocol it
replaces, which required a dedicated \texttt{scheduler\_lock} table, periodic
heartbeat renewal, and explicit expiration logic.

\begin{algorithm}[H]
\caption{Fair Job Queueing (Server-Side)}
\begin{algorithmic}[1]
\Procedure{QueueJobs}{$W_{max}, T_w$}

  \State \Comment{Advisory lock prevents concurrent replicas from over-allocating}
  \State $\text{locked} \gets D.\text{Query}(\texttt{SELECT pg\_try\_advisory\_lock(}
    \textit{queue\_lock\_id}\texttt{)})$
  \If{\textbf{not} $\text{locked}$}
    \State \Return
    \Comment{Another replica is already queueing; skip this cycle}
  \EndIf

  \State \Comment{--- All reads and writes below are inside the advisory lock ---}

  \State \Comment{Calculate available capacity}
  \State $W_{busy} \gets D.\text{CountWorkers}(\varsigma = \text{Busy})$
  \State $J_{queued} \gets D.\text{CountJobs}(\sigma = \text{Queued})$
  \State $W_a \gets W_{max} - W_{busy} - J_{queued}$
  \Comment{Slots not busy and not already queued}

  \If{$W_a \leq 0$}
    \State $D.\text{Query}(\texttt{SELECT pg\_advisory\_unlock(}\textit{queue\_lock\_id}\texttt{)})$
    \State \Return
  \EndIf

  \State $J_{pending} \gets D.\text{GetJobs}(\sigma = \text{Pending})$
  \If{$|J_{pending}| = 0$}
    \State $D.\text{Query}(\texttt{SELECT pg\_advisory\_unlock(}\textit{queue\_lock\_id}\texttt{)})$
    \State \Return
  \EndIf

  \State \Comment{Compute per-user fairness metrics}
  \State $U_{pending} \gets \{u_i : J_i \in J_{pending}\}$

  \ForAll{$u \in U_{pending}$}
    \State $n_u^r \gets D.\text{CountJobs}(u, \text{InProgress})$
    \State $n_u^p \gets D.\text{CountJobs}(u, \text{Pending})$
    \State $\tau_u \gets D.\text{SumRecentRuntime}(u, T_w)$
  \EndFor

  \State $\text{sorted} \gets \text{Sort}(U_{pending}, \text{key} = (n_u^r, \tau_u, -n_u^p))$
  \State $\text{allocations} \gets \text{AllocateSlots}(\text{sorted}, W_a)$
  \Comment{Round-robin, same as before}

  \State \Comment{Transition selected jobs to Queued}
  \ForAll{$u \in \text{sorted}$}
    \State $\text{jobs} \gets D.\text{GetJobs}(u, \text{Pending}, \text{limit} = \text{allocations}[u])$
    \ForAll{$J_i \in \text{jobs}$}
      \State $D.\text{UpdateJob}(J_i, \sigma_i = \text{Queued})$
    \EndFor
  \EndFor

  \State $D.\text{Query}(\texttt{SELECT pg\_advisory\_unlock(}\textit{queue\_lock\_id}\texttt{)})$

\EndProcedure
\end{algorithmic}
\end{algorithm}

\subsection{Step 2: Worker Claims Job (Atomic, No Leader)}

Workers claim the next available \texttt{Queued} job using a single atomic SQL
statement. PostgreSQL row-level locking (\texttt{FOR UPDATE SKIP LOCKED}) ensures
that concurrent workers never claim the same job.

\begin{algorithm}[H]
\caption{Atomic Job Claim (Worker-Side)}
\begin{algorithmic}[1]
\Procedure{ClaimJob}{$\omega_k$}

  \State \textbf{Execute atomic SQL:}
  \State \quad \texttt{UPDATE job\_info}
  \State \quad \texttt{SET status = 'InProgress',}
  \State \quad \quad \texttt{worker\_id = }$\omega_k$\texttt{,}
  \State \quad \quad \texttt{started\_at = NOW()}
  \State \quad \texttt{WHERE (session\_id, owner, repo) = (}
  \State \quad \quad \texttt{SELECT session\_id, owner, repo}
  \State \quad \quad \texttt{FROM job\_info}
  \State \quad \quad \texttt{WHERE status = 'Queued'}
  \State \quad \quad \texttt{ORDER BY created\_at}
  \State \quad \quad \texttt{LIMIT 1}
  \State \quad \quad \texttt{FOR UPDATE SKIP LOCKED}
  \State \quad \texttt{)}
  \State \quad \texttt{RETURNING *}

  \State \Return claimed job $J_i$ or $\emptyset$

\EndProcedure
\end{algorithmic}
\end{algorithm}

\textbf{Key property}: No leader election is needed. Postgres MVCC and
\texttt{SKIP LOCKED} provide serializable job claiming. Multiple workers calling
\texttt{ClaimJob} concurrently will each get a distinct job (or $\emptyset$ if none
are available).

\subsection{Fairness Preservation}

The fair scheduling priority $\pi(u) = (n_u^r, \tau_u, -n_u^p)$ is enforced in
Step~1 when the server selects which jobs to transition to \texttt{Queued}. Workers
in Step~2 then claim in FIFO order within the queued set. This preserves the
round-robin allocation from the current scheduler while removing the need for a
single leader to push jobs.


\section{Worker Pool Scaling}
\label{sec:scaling}

The pool manager (part of the server) adjusts $|P_\omega|$ based on workload.
Scaling uses Argo Workflows to create/delete worker pods, but workers are long-lived
and process multiple jobs.

\subsection{Scaling Policy}

\begin{equation}
W_{target} = \begin{cases}
  \min\bigl(W_{max},\; \max(\lambda \cdot W_{active},\; 1)\bigr)
    & \text{if } |J_{pending}| + |J_{queued}| > 0 \text{ and } W_{active} < W_{max} \\[6pt]
  \max\bigl(W_{min},\; W_{active} - \gamma\bigr)
    & \text{if } |J_{pending}| + |J_{queued}| = 0 \text{ and } t_{all\_idle} > T_{idle} \\[6pt]
  W_{active}
    & \text{otherwise}
\end{cases}
\end{equation}

where:
\begin{itemize}
\item $\lambda = 2$: ramp-up doubles workers each scaling interval (fast ramp)
\item $\gamma = 1$: decay removes one worker per decay interval (slow decay)
\item $t_{all\_idle}$: duration since all workers became idle
\end{itemize}

\subsection{Ramp-Up Sequence Example}

Starting from $W_{active} = 0$ with 100 pending jobs and $W_{max} = 16$:

\begin{center}
\begin{tabular}{cccc}
\toprule
\textbf{Interval} & $W_{active}$ & $W_{target}$ & \textbf{Workers Started} \\
\midrule
$t_0$ & 0 & 1 & 1 \\
$t_0 + T_{scale}$ & 1 & 2 & 1 \\
$t_0 + 2T_{scale}$ & 2 & 4 & 2 \\
$t_0 + 3T_{scale}$ & 4 & 8 & 4 \\
$t_0 + 4T_{scale}$ & 8 & 16 & 8 \\
\bottomrule
\end{tabular}
\end{center}

\subsection{Decay Sequence Example}

From $W_{active} = 10$, no pending jobs, $W_{min} = 0$, $\gamma = 1$:

After $T_{idle}$ elapses with no work, one worker is marked \texttt{Draining} every
$T_{decay}$ seconds. Draining workers finish their current job (if any) and exit.
Full drain: $10 \times T_{decay}$ seconds.


\subsection{Pool Manager Algorithm}

\begin{algorithm}[H]
\caption{Pool Scaling (Server-Side)}
\begin{algorithmic}[1]
\Procedure{AdjustPool}{$P_\omega, W_{min}, W_{max}, \lambda, \gamma$}

  \State \Comment{------ Reap Dead Workers ------}
  \ForAll{$\omega_k \in P_\omega$}
    \If{$t_{now} - \omega_k.\text{heartbeat\_at} > T_{hb}^{exp}$}
      \State $D.\text{UpdateWorker}(\omega_k, \varsigma_k = \text{Exited})$
      \State Requeue any job held by $\omega_k$: $\sigma_i \gets \text{Queued}$
    \EndIf
  \EndFor

  \State $W_{active} \gets D.\text{CountWorkers}(\varsigma \in \{\text{Idle}, \text{Busy}\})$

  \State \Comment{------ Scale Up ------}
  \State $J_{pending} \gets D.\text{CountJobs}(\sigma = \text{Pending})$
  \State $J_{queued} \gets D.\text{CountJobs}(\sigma = \text{Queued})$
  \State $\text{demand} \gets J_{pending} + J_{queued}$

  \If{$\text{demand} > 0$ \textbf{and} $W_{active} < W_{max}$}
    \State $W_{target} \gets \min(W_{max},\; \max(\lambda \cdot W_{active},\; 1))$
    \State $\Delta \gets W_{target} - W_{active}$
    \For{$i = 1$ \textbf{to} $\Delta$}
      \State $\omega_{new} \gets \Omega.\text{CreateWorkerPod}()$
      \Comment{Submit long-lived worker via Argo}
      \State $D.\text{RegisterWorker}(\omega_{new}, \varsigma = \text{Idle})$
    \EndFor
  \EndIf

  \State \Comment{------ Scale Down (slow decay) ------}
  \If{$\text{demand} = 0$ \textbf{and} $W_{active} > W_{min}$}
    \If{all workers idle for $> T_{idle}$}
      \State $n_{drain} \gets \min(\gamma, \; W_{active} - W_{min})$
      \State Select $n_{drain}$ idle workers
      \ForAll{selected $\omega_k$}
        \State $D.\text{UpdateWorker}(\omega_k, \varsigma_k = \text{Draining})$
        \Comment{Worker will exit on next poll}
      \EndFor
    \EndIf
  \EndIf

\EndProcedure
\end{algorithmic}
\end{algorithm}


\section{Worker Lifecycle}
\label{sec:worker-lifecycle}

Each worker $\omega_k$ runs a simple loop: register, claim, execute, repeat.

\subsection{Worker State Machine}

\begin{equation}
\varsigma_k: \text{Idle}
  \underset{\text{job completes}}{\overset{\text{claim job}}{\rightleftarrows}}
  \text{Busy}
  \qquad
  \text{Idle} \xrightarrow{T_{idle}\text{ or Draining}} \text{Exited}
\end{equation}

\begin{itemize}
\item \textbf{Idle}: Registered, polling for jobs. Sends heartbeats.
\item \textbf{Busy}: Executing a claimed job $J_i$. Sends heartbeats.
\item \textbf{Draining}: Pool manager has requested shutdown. Worker finishes current
  job (if any), then exits.
\item \textbf{Exited}: Worker has terminated. Pool manager may clean up.
\end{itemize}


\section{Full Round-Trip Representation}
\label{sec:full-round-trip}

The complete round-trip with the worker pool:

\begin{equation}
\begin{split}
C &\xrightarrow{(\mathcal{Q}, \mathcal{R}, u)} S \\
S &\xrightarrow{\forall \mathcal{R}_i: (J_i, u_i, \sigma_i=\text{Pending})} D \\
S &\xrightarrow[\text{fair schedule}]{\pi(u), W_a} D: \sigma_i \gets \text{Queued}
  \quad \text{(selected jobs)} \\
\omega_k &\xleftarrow[\text{claim}]{\texttt{FOR UPDATE SKIP LOCKED}} D:
  \sigma_i \gets \text{InProgress} \\
\omega_k &\xrightarrow{\text{execute}(\mathcal{Q}, \mathcal{R}_i)}
  \mathcal{R}_i^{\mathcal{Q}} \\
\omega_k &\xrightarrow{\text{store results}} \text{MinIO} \\
\omega_k &\xrightarrow{\sigma_i = \text{Succeeded}} D \\
\omega_k &\xleftarrow{\text{claim next}} D
  \quad \text{(worker stays alive)} \\
S &\xrightarrow{\mathcal{R}_i^{\mathcal{Q}}} C
\end{split}
\end{equation}

Detailed explanation:
\begin{enumerate}
\item $C \to S$: Client submits query suite $\mathcal{Q}$, repository list
  $\mathcal{R}$, and user $u$ to server.
\item $S \to D$: Server creates jobs $J_i$ for each repository with
  $\sigma_i = \text{Pending}$.
\item $S \to D$: Server runs fair scheduling, transitions selected jobs to
  $\sigma_i = \text{Queued}$.
\item $\omega_k \leftarrow D$: Worker atomically claims a queued job using
  \texttt{FOR UPDATE SKIP LOCKED}.
\item $\omega_k$: Worker executes queries and generates results.
\item $\omega_k \to \text{MinIO}$: Worker stores results.
\item $\omega_k \to D$: Worker marks job as Succeeded.
\item $\omega_k \leftarrow D$: Worker claims \emph{next} job (no exit/restart).
\item $S \to C$: Server forwards results to client.
\end{enumerate}


\section{Properties and Invariants}

\subsection{Capacity Constraint}

At all times:
\begin{equation}
|\{i : \sigma_i = \text{InProgress}\}| + |\{i : \sigma_i = \text{Queued}\}| \leq W_{max}
\end{equation}

This is enforced by the server's \textsc{QueueJobs} procedure.

\subsection{Claim Exclusivity}

No two workers claim the same job:
\begin{equation}
\forall J_i: |\{\omega_k : \omega_k.\text{current\_job} = J_i\}| \leq 1
\end{equation}

This is enforced by Postgres row-level locking (\texttt{FOR UPDATE SKIP LOCKED}).

\subsection{Worker Liveness}

A worker $\omega_k$ is considered alive if and only if:
\begin{equation}
t_{now} - \omega_k.\text{heartbeat\_at} < T_{hb}^{exp}
\end{equation}

Dead workers (violating this) have their jobs requeued by the pool manager.

\subsection{Scaling Bounds}

The pool size is always bounded:
\begin{equation}
W_{min} \leq |P_\omega| \leq W_{max}
\end{equation}

\subsection{Fairness Preservation}

The priority ordering $\pi(u) = (n_u^r, \tau_u, -n_u^p)$ is applied identically to
the current scheduler. The only change is that selected jobs transition to
\texttt{Queued} instead of being submitted to Argo directly.


\section{Configuration Parameters}

\begin{center}
\begin{longtable}{p{3.5cm}p{1.5cm}p{2cm}p{7cm}}
\toprule
\textbf{Parameter} & \textbf{Symbol} & \textbf{Default} & \textbf{Description} \\
\midrule
\endfirsthead

\toprule
\textbf{Parameter} & \textbf{Symbol} & \textbf{Default} & \textbf{Description} \\
\midrule
\endhead

\midrule
\endfoot

\bottomrule
\endlastfoot

Min Workers & $W_{min}$ & 0 & Minimum pool size; set $> 0$ for always-warm pool \\
Max Workers & $W_{max}$ & 10 & Maximum concurrent workers \\
Ramp Factor & $\lambda$ & 2 & Multiplicative scale-up factor \\
Decay Rate & $\gamma$ & 1 & Workers removed per decay step \\
Idle Timeout & $T_{idle}$ & 300s & Idle duration before worker becomes eligible for
  removal \\
Heartbeat Interval & $T_{hb}$ & 10s & Worker heartbeat frequency \\
Heartbeat Expiry & $T_{hb}^{exp}$ & 30s & Time after which a silent worker is
  considered dead \\
Scale Check & $T_{scale}$ & 15s & Pool size evaluation frequency \\
Decay Interval & $T_{decay}$ & 60s & Minimum time between scale-down steps \\
Claim Poll & $T_{claim}$ & 2s & Worker job-claim polling frequency \\
Runtime Window & $T_w$ & 3600s & Window for $\tau_u$ calculation \\
\end{longtable}
\end{center}

Environment variables (proposed):
\begin{itemize}
\item \texttt{MRVA\_POOL\_MIN\_WORKERS} -- Sets $W_{min}$
\item \texttt{MRVA\_POOL\_MAX\_WORKERS} -- Sets $W_{max}$
\item \texttt{MRVA\_POOL\_RAMP\_FACTOR} -- Sets $\lambda$
\item \texttt{MRVA\_POOL\_DECAY\_RATE} -- Sets $\gamma$
\item \texttt{MRVA\_POOL\_IDLE\_TIMEOUT\_SECONDS} -- Sets $T_{idle}$
\item \texttt{MRVA\_POOL\_HEARTBEAT\_SECONDS} -- Sets $T_{hb}$
\item \texttt{MRVA\_POOL\_SCALE\_CHECK\_SECONDS} -- Sets $T_{scale}$
\item \texttt{MRVA\_POOL\_CLAIM\_POLL\_SECONDS} -- Sets $T_{claim}$
\end{itemize}


\newpage
\section{Execution Loop in Pseudo-Code}

\noindent\textbf{Listing: Server-Side Pool Manager and Fair Queueing}

\begin{lstlisting}[language=Python]
# Server-side: Pool manager and fair job queueing
# Replaces per-job Argo submission with pool scaling + job queueing

# Initialization (*@\gomark{pkg/pool/pool.go}{26}@*)
$D$ = PostgresDB()
$S$ = Server()
$\Omega$ = ArgoWorkflows()  # Used to create/delete worker pods
$P_\omega$ = set()  # Worker pool (tracked in Postgres)
$W_{min}$ = 0
$W_{max}$ = 10
$\lambda$ = 2  # Ramp factor
$\gamma$ = 1  # Decay rate
$T_w$ = 3600  # Runtime window

# Client submits queries (*@\gomark{pkg/server/server.go}{1019}@*)
$C$ submits ($\mathcal{Q}$, $\mathcal{R}$, $u$) to $S$

# Server stores jobs in database (*@\gomark{pkg/server/server.go}{26}@*)
for $\mathcal{R}_i$ in $\mathcal{R}$:
    $J_i$ = JobSpec(session_id, owner, repo)
    $D$.store_job($J_i$, user=$u$, status="Pending", payload=$\mathcal{Q}$)

# Server main loop (no leader election needed) (*@\gomark{pkg/pool/pool.go}{120}@*)
while True:
    # --- Reap dead workers --- (*@\gomark{pkg/pool/pool.go}{342}@*)
    for $\omega_k$ in $D$.get_workers():
        if $t_{now}$ - $\omega_k$.heartbeat_at > $T_{hb}^{exp}$:
            $D$.update_worker($\omega_k$, status="Exited")
            if $\omega_k$.current_job is not None:
                $D$.update_job($\omega_k$.current_job, status="Queued",
                              worker_id=None)

    # --- Fair scheduling: Pending $\to$ Queued --- (*@\gomark{pkg/pool/pool.go}{210}@*)
    # Advisory lock prevents concurrent replicas from over-allocating
    if $D$.query("SELECT pg_try_advisory_lock(queue_lock_id)"):

        $W_{busy}$ = $D$.count_workers(status="Busy")
        $J_{queued}$ = $D$.count_jobs(status="Queued")
        $W_a$ = $W_{max}$ - $W_{busy}$ - $J_{queued}$

        if $W_a$ > 0:
            $J_{pending}$ = $D$.get_jobs(status="Pending")
            $U_{pending}$ = {$u_i$ for $J_i$ in $J_{pending}$}

            # Calculate fairness metrics
            for $u$ in $U_{pending}$:
                $n_u^r$ = $D$.count_jobs(user=$u$, status="InProgress")
                $n_u^p$ = $D$.count_jobs(user=$u$, status="Pending")
                $\tau_u$ = $D$.sum_recent_runtime(user=$u$, window=$T_w$)

            sorted_users = sorted($U_{pending}$, (*@\gomark{pkg/pool/pool.go}{541}@*)
                                  key=lambda u: ($n_u^r$, $\tau_u$, -$n_u^p$))
            allocations = allocate_slots(sorted_users, $W_a$)

            for $u$ in sorted_users:
                jobs = $D$.get_jobs(user=$u$, status="Pending",
                                  limit=allocations[$u$])
                for $J_i$ in jobs:
                    $D$.update_job($J_i$, status="Queued")

            # --- Submit queued jobs to Argo --- (*@\gomark{pkg/pool/pool.go}{309}@*)
            # Each job spawns a worker pod (no --job flag in pool mode);
            # the worker claims jobs from Postgres via ClaimJob.
            for $J_i$ in newly_queued:
                uid = $\Omega$.submit_workflow($J_i$) (*@\gomark{pkg/scheduler/argo\_client.go}{502}@*)
                $D$.update_job_info($J_i$, argo_workflow_id=uid,
                                   started_at=$t_{now}$)

        $D$.query("SELECT pg_advisory_unlock(queue_lock_id)")

    # --- Pool scaling (implicit) ---
    # Worker creation is integrated into Argo submission above:
    # one workflow per queued job, capped by $W_a$.
    # Workers self-exit on idle timeout; no explicit scale-down step.

    # --- Reconcile Postgres with Argo ground truth --- (*@\gomark{pkg/pool/pool.go}{370}@*)
    # Advisory-locked; only one replica polls at a time
    if $D$.query("SELECT pg_try_advisory_lock(argo_lock_id)"):
        for $J_i$ in $D$.get_jobs(status in ["InProgress", "Queued"]):
            if $J_i$.argo_workflow_id is not None:
                phase = $\Omega$.get_workflow_status($J_i$.argo_workflow_id)
                if phase in ["Succeeded", "Failed", "Error"]:
                    $D$.update_job($J_i$, status=map_phase(phase),
                                  completed_at=$t_{now}$)
        $D$.query("SELECT pg_advisory_unlock(argo_lock_id)")

    # --- Result delivery (separate goroutine) --- (*@\gomark{pkg/server/server.go}{1144}@*)
    # Agents publish results to RabbitMQ after CompleteJob;
    # server reads from queue, stores to state, triggers scheduling.

    sleep($T_{scale}$)
\end{lstlisting}

\FloatBarrier
\newpage

\noindent\textbf{Listing: Worker Execution Loop (Inside Long-Lived Pod)}

\begin{lstlisting}[language=Python]
# Worker execution loop inside a long-lived Argo pod
# Worker persists across jobs: claim $\to$ execute $\to$ complete $\to$ claim next

# Initialization
$\omega_k$ = generate_worker_id()
$D$ = PostgresDB()
MinIO = MinIOClient(env["MRVA_BCS_ENDPOINT"], env["MRVA_BUCKET"], ...)
HEPC = HEPCClient(env["MRVA_HEPC_ENDPOINT"])
codeql_cli = env["CODEQL_CLI_PATH"]

# Register with pool (*@\gomark{pkg/state/state\_postgres.go}{938}@*)
$D$.register_worker($\omega_k$, status="Idle")
$t_{idle\_start}$ = $t_{now}$

# Main worker loop
while True:
    # Check if pool manager wants us to drain (*@\gomark{pkg/state/state\_postgres.go}{978}@*)
    $\varsigma_k$ = $D$.get_worker_status($\omega_k$)
    if $\varsigma_k$ == "Draining":
        $D$.update_worker($\omega_k$, status="Exited")
        exit(0)

    # Heartbeat (*@\gomark{pkg/state/state\_postgres.go}{959}@*)
    $D$.update_worker($\omega_k$, heartbeat_at=$t_{now}$)

    # Attempt to claim a job (atomic, via Postgres SKIP LOCKED) (*@\gomark{pkg/state/state\_postgres.go}{1027}@*)
    $J_i$ = claim_job($\omega_k$)  # Returns None if no queued jobs

    if $J_i$ is None:
        # No work available
        if $t_{now}$ - $t_{idle\_start}$ > $T_{idle}$:
            # Self-exit after idle timeout
            $D$.update_worker($\omega_k$, status="Exited")
            exit(0)
        sleep($T_{claim}$)
        continue

    # --- Got a job: execute it --- (*@\gomark{pkg/agent/agent.go}{88}@*)
    $t_{idle\_start}$ = $t_{now}$  # Reset idle timer
    $D$.update_worker($\omega_k$, status="Busy", current_job=$J_i$)

    # Fetch query pack from MinIO (*@\gomark{pkg/agent/agent.go}{105}@*)
    $\mathcal{Q}$ = MinIO.get($J_i$.query_pack_location)
    extract_query_pack($\mathcal{Q}$, "/tmp/queries")

    # Fetch CodeQL database from HEPC (*@\gomark{pkg/agent/agent.go}{125}@*)
    db_url = HEPC.get_database_url($J_i$.owner, $J_i$.repo, $J_i$.language)
    $DB_i$ = download_and_extract(db_url, "/tmp/databases")

    # Execute CodeQL queries (*@\gomark{pkg/codeql/codeql.go}{126}@*)
    $\mathcal{R}_i^{\mathcal{Q}}$ = set()
    for $\mathcal{Q}_j$ in $\mathcal{Q}$:
        results = run_command([
            codeql_cli, "database", "analyze",
            $DB_i$, "/tmp/queries",
            "--format=sarif-latest",
            "--output=/tmp/results.sarif"
        ])
        for result in parse_sarif("/tmp/results.sarif"):
            $\mathcal{R}_i^{\mathcal{Q}}$.add(result)

    # Store results to MinIO (*@\gomark{pkg/agent/agent.go}{159}@*)
    results_loc = f"{$J_i$.session_id}/{$J_i$.owner}/{$J_i$.repo}/results.json"
    MinIO.put(results_loc, $\mathcal{R}_i^{\mathcal{Q}}$)

    # Mark job complete (Postgres) (*@\gomark{pkg/state/state\_postgres.go}{1096}@*)
    $D$.update_job($J_i$, status="Succeeded", completed_at=$t_{now}$)
    $D$.update_worker($\omega_k$, status="Idle", current_job=None,
                     jobs_completed=$\omega_k$.jobs_completed + 1)

    # Publish result to RabbitMQ for server ConsumeResults()
    RabbitMQ.publish_result($J_i$, $\mathcal{R}_i^{\mathcal{Q}}$)

    # Continue loop: claim next job (no exit!)
\end{lstlisting}

\FloatBarrier
\newpage

\noindent\textbf{Listing: Atomic Job Claim via Postgres}
\gomark{pkg/state/state\_postgres.go}{1027}

\begin{lstlisting}[language=SQL, numbers=none, mathescape=false]
-- Worker claims next available queued job (inside a transaction).
-- FOR UPDATE SKIP LOCKED ensures no double-claiming.
-- Multiple workers calling concurrently each get a distinct job.

BEGIN;

-- Step 1: atomically select and lock a queued job
SELECT session_id, owner, repo
FROM job_status
WHERE status = 'Queued'
ORDER BY session_id, owner, repo
LIMIT 1
FOR UPDATE SKIP LOCKED;

-- Step 2: transition the job to InProgress
UPDATE job_status
SET status = 'InProgress'
WHERE session_id = $1 AND owner = $2 AND repo = $3;

-- Step 3: record worker assignment and start time
UPDATE job_info
SET worker_id = $4, started_at = NOW()
WHERE session_id = $1 AND owner = $2 AND repo = $3;

COMMIT;
\end{lstlisting}


\newpage
\section{Execution Loop in Pseudo-Code, algorithmic}

\noindent\textbf{Algorithm: Server-side pool management and fair job queueing}

\begin{algorithmic}[1]
    \Procedure{ServerPoolLoop}{$W_{min}, W_{max}, \lambda, \gamma, T_w$}

      \State \Comment{------ Initialization ------}
      \gomark{pkg/pool/pool.go}{26}
      \State $D \gets$ PostgreSQL database
      \State $S \gets$ Server (API coordinator)
      \State $\Omega \gets$ Argo Workflows (for pod lifecycle only)
      \State MinIO $\gets$ Object storage
      \State HEPC $\gets$ CodeQL database server

      \While{\textbf{not} stopped}

        \State \Comment{====== Reap Dead Workers ======}
        \gomark{pkg/pool/pool.go}{342}
        \ForAll{$\omega_k : t_{now} - \omega_k.\text{heartbeat\_at} > T_{hb}^{exp}$}
          \State $D.\text{UpdateWorker}(\omega_k, \varsigma_k = \text{Exited})$
          \If{$\omega_k.\text{current\_job} = J_i \neq \text{NULL}$}
            \State $D.\text{UpdateJob}(J_i, \sigma_i = \text{Queued}, \text{worker\_id} = \text{NULL})$
            \Comment{Requeue}
          \EndIf
        \EndFor

        \State \Comment{====== Fair Scheduling: Pending $\to$ Queued ======}
        \gomark{pkg/pool/pool.go}{210}
        \State $\text{locked} \gets D.\texttt{pg\_try\_advisory\_lock}(\textit{queue\_lock\_id})$
        \Comment{Prevent concurrent over-allocation}
        \If{$\text{locked}$}
          \State $W_{busy} \gets D.\text{CountWorkers}(\varsigma = \text{Busy})$
          \State $J_{queued\_count} \gets D.\text{CountJobs}(\sigma = \text{Queued})$
          \State $W_a \gets W_{max} - W_{busy} - J_{queued\_count}$

          \If{$W_a > 0$}
            \State $J_{pending} \gets D.\text{GetJobs}(\sigma = \text{Pending})$
            \State $U_{pending} \gets \{u_i : J_i \in J_{pending}\}$

            \ForAll{$u \in U_{pending}$}
              \State $n_u^r \gets D.\text{CountJobs}(u, \text{InProgress})$
              \State $n_u^p \gets D.\text{CountJobs}(u, \text{Pending})$
              \State $\tau_u \gets D.\text{SumRecentRuntime}(u, T_w)$
            \EndFor

            \State $\text{sorted} \gets \text{Sort}(U_{pending}, \text{key} = (n_u^r, \tau_u, -n_u^p))$
            \gomark{pkg/pool/pool.go}{541}
            \State $\text{alloc} \gets \text{AllocateSlots}(\text{sorted}, W_a)$
            \Comment{Round-robin}

            \ForAll{$u \in \text{sorted}$}
              \State $\text{jobs} \gets D.\text{GetJobs}(u, \text{Pending}, \text{limit} = \text{alloc}[u])$
              \ForAll{$J_i \in \text{jobs}$}
                \State $D.\text{UpdateJob}(J_i, \sigma_i = \text{Queued})$
                \Comment{Available for workers}
              \EndFor
            \EndFor

            \State \Comment{------ Submit to Argo ------}
            \gomark{pkg/pool/pool.go}{309}
            \ForAll{$J_i \in \text{newly queued}$}
              \State $\text{uid} \gets \Omega.\text{SubmitWorkflow}(J_i)$
              \gomark{pkg/scheduler/argo\_client.go}{502}
              \Comment{Pool mode: no \texttt{--job} flag}
              \State $D.\text{SetJobInfo}(J_i, \text{argo\_workflow\_id} = \text{uid},
                \text{started\_at} = t_{now})$
            \EndFor
          \EndIf
          \State $D.\texttt{pg\_advisory\_unlock}(\textit{queue\_lock\_id})$
        \EndIf

        \State \Comment{====== Pool Scaling (implicit) ======}
        \State \Comment{Worker creation is integrated into Argo submission above;}
        \State \Comment{one workflow per queued job, capped by $W_a$.}
        \State \Comment{Workers self-exit on idle timeout.}

        \State \Comment{====== Reconcile with Argo Ground Truth ======}
        \gomark{pkg/pool/pool.go}{370}
        \State $\text{locked} \gets D.\texttt{pg\_try\_advisory\_lock}(\textit{argo\_lock\_id})$
        \If{$\text{locked}$}
          \ForAll{$J_i : \sigma_i \in \{\text{InProgress}, \text{Queued}\}$
            and $J_i.\text{argo\_workflow\_id} \neq \text{NULL}$}
            \State $\text{phase} \gets \Omega.\text{GetWorkflowStatus}(J_i.\text{argo\_workflow\_id})$
            \If{$\text{phase} \in \{\text{Succeeded}, \text{Failed}, \text{Error}\}$}
              \State $D.\text{UpdateJob}(J_i, \sigma_i = \text{MapPhase}(\text{phase}),
                t_i^c = t_{now})$
            \EndIf
          \EndFor
          \State $D.\texttt{pg\_advisory\_unlock}(\textit{argo\_lock\_id})$
        \EndIf

        \State \Comment{====== Result Delivery (separate goroutine) ======}
        \gomark{pkg/server/server.go}{1144}
        \State \Comment{Agents publish results to RabbitMQ after CompleteJob;}
        \State \Comment{server reads from queue, stores to state, triggers scheduling.}

        \State Sleep($T_{scale}$)
      \EndWhile

    \EndProcedure
\end{algorithmic}

\vspace{1em}

\noindent\textbf{Algorithm: Worker execution loop inside long-lived pod}

\begin{algorithmic}[1]
    \Procedure{WorkerLoop}{$\omega_k$, env\_vars}

      \State \Comment{------ Initialization ------}
      \State $D \gets$ PostgreSQL database
      \State MinIO $\gets$ \text{MinIOClient}(\text{env})
      \State HEPC $\gets$ \text{HEPCClient}(\text{env})
      \State $D.\text{RegisterWorker}(\omega_k, \varsigma_k = \text{Idle})$
      \gomark{pkg/state/state\_postgres.go}{938}
      \State $t_{idle\_start} \gets t_{now}$

      \While{\textbf{true}}

        \State \Comment{------ Check for Drain Signal ------}
        \gomark{pkg/state/state\_postgres.go}{978}
        \State $\varsigma_k \gets D.\text{GetWorkerStatus}(\omega_k)$
        \If{$\varsigma_k = \text{Draining}$}
          \State $D.\text{UpdateWorker}(\omega_k, \varsigma_k = \text{Exited})$
          \State \textbf{exit}(0)
        \EndIf

        \State \Comment{------ Heartbeat ------}
        \gomark{pkg/state/state\_postgres.go}{959}
        \State $D.\text{Heartbeat}(\omega_k, t_{now})$

        \State \Comment{------ Claim Job (atomic via Postgres MVCC) ------}
        \gomark{pkg/state/state\_postgres.go}{1027}
        \State $J_i \gets \text{ClaimJob}(\omega_k)$
        \Comment{\texttt{FOR UPDATE SKIP LOCKED}}

        \If{$J_i = \emptyset$}
          \If{$t_{now} - t_{idle\_start} > T_{idle}$}
            \State $D.\text{UpdateWorker}(\omega_k, \varsigma_k = \text{Exited})$
            \State \textbf{exit}(0) \Comment{Self-exit after idle timeout}
          \EndIf
          \State Sleep($T_{claim}$)
          \State \textbf{continue}
        \EndIf

        \State \Comment{------ Execute Job ------}
        \gomark{pkg/agent/agent.go}{88}
        \State $t_{idle\_start} \gets t_{now}$ \Comment{Reset idle timer}
        \State $D.\text{UpdateWorker}(\omega_k, \varsigma_k = \text{Busy},
          \text{current\_job} = J_i)$

        \State $\mathcal{Q} \gets \text{MinIO}.\text{Get}(J_i.\text{query\_pack\_location})$
        \Comment{Fetch query pack}
        \State $DB_i \gets \text{HEPC}.\text{FetchDatabase}(J_i.\text{owner},
          J_i.\text{repo})$
        \Comment{Fetch CodeQL DB}

        \State $\mathcal{R}_i^{\mathcal{Q}} \gets \emptyset$
        \ForAll{$\mathcal{Q}_j \in \mathcal{Q}$}
          \State Run CodeQL CLI: \texttt{database analyze} $DB_i$ with $\mathcal{Q}_j$
          \gomark{pkg/codeql/codeql.go}{126}
          \State $\mathcal{R}_i^{\mathcal{Q}} \gets \mathcal{R}_i^{\mathcal{Q}} \cup
            \text{results}$
        \EndFor

        \State \Comment{------ Store Results and Complete ------}
        \gomark{pkg/state/state\_postgres.go}{1096}
        \State $\text{MinIO}.\text{Put}(\text{BuildPath}(J_i), \mathcal{R}_i^{\mathcal{Q}})$
        \State $D.\text{UpdateJob}(J_i, \sigma_i = \text{Succeeded}, t_i^c = t_{now})$
        \State $D.\text{UpdateWorker}(\omega_k, \varsigma_k = \text{Idle},
          \text{current\_job} = \text{NULL})$
        \State $\text{RabbitMQ}.\text{PublishResult}(J_i, \mathcal{R}_i^{\mathcal{Q}})$
        \Comment{For server ConsumeResults}

        \State \Comment{Loop continues: claim next job without exit}
      \EndWhile

    \EndProcedure
\end{algorithmic}


\section{Implementation Notes}

\subsection{Code Locations (Proposed)}

\begin{itemize}
\item \textbf{Shared (mrvacommander/)}:
  \begin{itemize}
  \item Pool interface types and Postgres queries:
    \texttt{pkg/pool/pool.go}
  \item Job claim SQL (\texttt{FOR UPDATE SKIP LOCKED}):
    \texttt{pkg/state/state\_postgres.go} (extend)
  \item Worker registration and heartbeat:
    \texttt{pkg/state/state\_postgres.go} (extend)
  \item Database migration:
    \texttt{db\_schema\_changes/pool/}
  \end{itemize}

\item \textbf{Agent (agent-commander-argo/)}:
  \begin{itemize}
  \item Worker loop (claim $\to$ execute $\to$ complete $\to$ claim):
    replace current single-job entrypoint
  \item Heartbeat goroutine
  \item Drain signal handling
  \end{itemize}

\item \textbf{Server (server-commander/)}:
  \begin{itemize}
  \item Pool manager goroutine: scaling decisions, worker reaping
  \item Fair queueing (Pending $\to$ Queued): adapt from current scheduler
  \item Result delivery loop: unchanged
  \end{itemize}
\end{itemize}

\subsection{Migration Path}

The worker pool replaces the existing scheduler and per-job pods outright.
Keeping the old paths alongside would risk jobs being processed by both systems
or, worse, silently working by accident and masking integration issues.

\begin{enumerate}
\item Add \texttt{worker\_id} column and \texttt{Queued} status to \texttt{job\_info}.
\item Replace the existing leader-elected scheduler with the pool manager
  (fair queueing via advisory lock, pool scaling). Remove the
  \texttt{scheduler\_lock} table and leader election code.
\item Replace the single-job agent entrypoint with the worker loop
  (claim $\to$ execute $\to$ complete $\to$ claim). Remove the per-job Argo
  workflow submission path (\texttt{SubmitWorkflow} per $J_i$).
\item Remove the Argo workflow status polling loop (workers now mark their own
  jobs as Succeeded/Failed directly in Postgres).
\end{enumerate}


\end{document}

%%% Local Variables:
%%% mode: LaTeX
%%% TeX-master: t
%%% TeX-engine: luatex
%%% TeX-command-extra-options: "-synctex=1 -shell-escape -interaction=nonstopmode"
%%% End:
